\chapter{Conclusion and Future Work}
\label{ch:conclusion}

\section{Summary}
\label{sec:c-summary}

This thesis asked how much language model reliable on-device action requires, and
answered it through A.L.F.R.E.D., a two-tier local-first middleware that routes
each request to a tiny reflexer or a heavier thinker. The five research questions
are settled by the evaluation.

\emph{On scale versus distillation} (RQ1), a 2B reflexer with a distilled pattern
matches a 35B free-form thinker on binding (both 100\% balanced), while the same
2B without a pattern reaches 30\%; the pattern is worth roughly $17\times$ the
parameters. \emph{On the spectrum} (RQ2), the abstraction is an efficiency win on
a clean API (a $+65$-point oracle gain at two calls instead of nine) and an
accuracy win on a realistic one. \emph{On the limit of scale} (RQ3), a 35B thinker
rescues discovery but not convention, failing device encodings consistently, while
a distilled pattern supplies exactly that convention. \emph{On teacher quality}
(RQ4), end-to-end student accuracy tracks the teacher: a cloud teacher's patterns
score 80\% against a local teacher's 60\% on the same reflexer and tasks. \emph{On
coverage and routing} (RQ5), the reflexer's territory is the deterministic head
that on-device assistants target, and the open bottleneck is routing, not binding.

The contribution, stated once more in plain terms: \textbf{on a real device API,
scale buys discovery but not undiscoverable convention; distillation buys the
convention; and a small reflexer applies it at a fraction of the cost.} The
executioner paradigm is validated through the full teacher-to-student loop, and
its central mechanism, compressing a teacher's device knowledge into an
inspectable symbolic pattern rather than into student weights, is the thesis's
main idea.

\section{Future work}
\label{sec:c-future}

The limitations of Chapter~\ref{ch:discussion} map onto five concrete next steps,
ordered by leverage.

\paragraph{Session-based device grounding.} Both teachers miss the genuinely
undiscoverable conventions because cold-start distillation gives them no way to
check a value against reality. Granting the teacher read-only access to the real
device, so it observes actual codes before committing an encoding, would convert
those guesses into facts and close the remaining convention errors. The integrity
rule carries over: ground on the device or a development split, never on the
evaluation gold.

\paragraph{Deterministic value transforms.} The reflexer's residual errors are
arithmetic and enum-lookup mistakes (percent-to-0--255, label-to-code). Moving
these transforms out of the model and into a small deterministic function inside
the pattern removes them for both the teacher and the student, and is a bounded
engineering change.

\paragraph{The router.} The evaluation localizes the open bottleneck precisely:
retrieval places the correct candidate in the top-3, but selection picks the wrong
sibling, and selection does not improve with selector scale. The next investment
is sibling discrimination, whether through better selection prompting or a learned
selector over behavioural signals, since the router gates both tiers.

\paragraph{Iterative, grounded distillation.} The cold-start patterns can be run,
executed on a real device, and re-distilled from the observed failures over two to
three rounds, lifting a weak local teacher toward cloud quality without a larger
model. This iteration lives entirely in the offline learning phase; the run-time
reflex remains a single pass over a frozen pattern, so the simplicity that makes
the reflexer deployable is preserved by construction, with a minimality check each
round to stop patterns accreting special cases.

\paragraph{Authoring and personalization at scale.} Longer-term, mining a user's
own repeated trajectories into candidate patterns would extend coverage beyond
hand-authored skills, and a privacy-preserving sharing layer would let patterns be
exchanged without leaking personal data. A parametric alternative, fine-tuning a
very small model on the pattern conventions, is a path worth measuring against the
symbolic approach, trading inspectability for a smaller artifact.

\paragraph{Development substrate.} All of the above extensions are intended to be
built inside the \textbf{Pi coding-agent framework}~\citep{pi-agent} on which the
thinker already runs (Section~\ref{sec:thinker}). Consolidating the reflexer's
future work in the same minimal, transparent environment keeps the executor and
the agent loop on one runtime, tool surface, and instrumentation path, which lowers
the cost of the iterate-and-measure loop the rest of this section depends on.

\section{Closing}
\label{sec:c-closing}

The shift toward on-device assistants raises a practical question that this thesis
takes seriously: not how powerful a model can be, but how little model an action
actually needs once its knowledge is placed where it belongs. The answer found
here is that a great deal of on-device action is not a reasoning problem at all. It
is a knowledge-placement problem, and once the convention is distilled into a
pattern, a small model executing a frozen reflex is enough.
